\documentclass[a4paper,11pt]{article}

\usepackage[ngerman]{babel}
\usepackage[top=3cm,bottom=3cm,left=3cm,right=3cm]{geometry} %NICHT IN KEMIE2


%\usepackage{microtype} %schöneres typesetting  mit [stretch=10] für nur 1% anstelle von 2%
%\usepackage{lmodern} %Latin Modern FONT
%\usepackage{pifont} %schönere Font



\usepackage{amsmath}
\usepackage{nicefrac}
\usepackage{amssymb}
%\usepackage{mathtools}

\usepackage{titlesec} %Um Section, etc bearbeiten zu können
\usepackage{setspace}
\usepackage{enumitem} %enumerate
\usepackage{multicol}
\usepackage[many]{tcolorbox}
\usepackage{hyperref}
\usepackage{arydshln}

%\renewcommand{\ttdefault}{lmtt} %Schriftart wird geändert zu Latin Modern Typewriter


\hypersetup{hidelinks,
			colorlinks=true,
			linkcolor=black,
			citecolor=miudiutex,
			urlcolor=miugray2
			}

\setlength{\parindent}{0pt}
\setcounter{section}{0}


\titlespacing*{\section}{0pt}{20mm}{6mm}
\titlespacing*{\subsection}{0pt}{6mm}{3mm}
\titlespacing*{\subsubsection}{0pt}{3mm}{0mm}



%\definecolor{miudiutex}{HTML}{2f3d39}
\definecolor{miudiutex}{HTML}{1d2926}
\definecolor{miugray}{HTML}{b4cccf}
\definecolor{miugray2}{HTML}{4d7365}
\definecolor{red}{HTML}{cf1f5c}
\definecolor{green}{HTML}{00A36C}
\definecolor{delphinidin}{HTML}{315794}
\definecolor{pelargonidin}{HTML}{b33807}
\definecolor{cyanidin}{HTML}{ba0746}
\definecolor{petunidin}{HTML}{b56fed}



\raggedcolumns %wichtig für Multicolumns


%================================
% SYMBOLS
%================================

\newcommand{\RA}{$\rightarrow$~}
\newcommand{\RL}{$\rightleftharpoons$~}
\newcommand{\HARP}{\ding{226}~}
%$\xrightarrow{2}$ %Pfeil mit 2 drüber


%================================
% SHORT COMMANDS (BOXES ; ETC)
%================================


\newcommand{\notiz}[1]{\begin{notizz}{}{}\begin{center}
			{\color{miugray2!50!miudiutex}#1} \end{center}\end{notizz}}
\newcommand{\zit}[2]{\begin{zitat*}{#1}{}\small #2 \end{zitat*}}

\newcommand{\frage}[2]{\begin{qst*}{#1}{}\small #2 \end{qst*}}

\newcommand{\unten}{\end{center} \tcblower \begin{center}}


\newcommand*\circled[1]{\tikz[baseline=(char.base)]{
            \node[shape=circle,draw,inner sep=0.5pt,miudiutex] (char) {\tiny #1};}}
            
\newcommand{\miusquare}{{\color{miugray2!50}\scriptsize $\blacksquare$}~}
\newcommand{\miusquarp}{{\color{petunidin!50}\scriptsize $\blacktriangleright$}~}
\newcommand{\niu}{\item[\miusquare]}
\newcommand{\nip}{\item[\miusquarp]}

%%%%% ABSAETZE

\newcommand{\pps}{\par \vspace*{1mm}}
\newcommand{\pp}{\par \vspace*{3mm}}
\newcommand{\ppbig}{\par \vspace*{8mm}}
\newcommand{\pphuge}{\par \vspace*{10mm}}

%================================
% FRAGE BOX
%================================

\tcbuselibrary{theorems,skins,hooks}
\newtcbtheorem{qst}{Frage:}
{%
	enhanced,
	breakable,
	colback = gray!15!white,
	frame hidden,
	boxrule = 0sp,
	borderline west = {2.5pt}{0pt}{red!70},
	sharp corners,
	detach title,
	before upper = \tcbtitle\par\smallskip,
	coltitle = black,
	fonttitle = \bfseries\sffamily,
	description font = \mdseries,
	separator sign none,
	segmentation style={solid, gray!50!black},
}
{th}




%================================
% ZITAT BOX
%================================

\tcbuselibrary{theorems,skins,hooks}
\newtcbtheorem{zitat}{Zitat}
{%
	enhanced,
	breakable,
	colback = gray!15!white,
	frame hidden,
	boxrule = 0sp,
	borderline west = {2.5pt}{0pt}{black!70},
	sharp corners,
	detach title,
%	before upper = \tcbtitle\par\smallskip,
	coltitle = gray!50!black,
	fonttitle = \bfseries\sffamily,
	description font = \mdseries,
%	separator sign none,
	segmentation style={solid, gray!50!black},
}
{th}



%================================
% SIMPLE SCHWARZER RAHMEN BOX
%================================


\newcommand{\boxfr}[1]{
	\begin{tcolorbox}[
	drop shadow southeast,
	enhanced,
	colframe=black!50,
	colback=white,
	boxrule = 1pt, 
%	toprule=0pt, bottomrule=0pt,rightrule=0pt,leftrule=0pt,
	boxsep=5pt,
	arc=1pt,
	]
	\begin{center}
	#1
	\end{center}
	\end{tcolorbox}
}

\definecolor{miudiutex}{HTML}{1d2926}
\definecolor{miugray}{HTML}{b4cccf}
\definecolor{miugray2}{HTML}{4d7365}
\definecolor{white2}{HTML}{f5f7f8}
\definecolor{red}{HTML}{cf1f5c}
\definecolor{green}{HTML}{00A36C}
\definecolor{delphinidin}{HTML}{315794}
\definecolor{uniblau}{HTML}{004e9f}
\definecolor{unigelb}{HTML}{fcba00}
\definecolor{unigrau}{HTML}{909085}
\definecolor{pelargonidin}{HTML}{b33807}
\definecolor{cyanidin}{HTML}{ba0746}
\definecolor{paeonidin}{HTML}{d15ca6}
\definecolor{petunidin}{HTML}{b56fed}
\definecolor{malvidin}{HTML}{8a2d8c}

\newcommand{\metermilli}{\mskip3mu \text{mm}}
\newcommand{\cm}{\mskip3mu \text{cm}}
\newcommand{\m}{ \text{m}}
\newcommand{\lite}{ \text{l}}
\newcommand{\kg}{ \text{kg}}
\newcommand{\g}{ \text{g}}
\newcommand{\km}{ \text{km}}
\newcommand{\h}{ \text{h}}
\newcommand{\s}{ \text{s}}
\newcommand{\tonne}{ \text{t}}
\usepackage[utf8]{inputenc}
\usepackage[T1]{fontenc}
\usepackage{lmodern}

\usepackage[
    activate={true,nocompatibility},
    final,
    tracking=true,
    kerning=true,
    spacing=true,
    factor=1100,
    stretch=30,
    shrink=20
]{microtype}
\usepackage{pdfpages}
\usepackage{awesomebox}
\usepackage{marginnote}
\tcbuselibrary{documentation}
\usepackage{sidenotes}
\usepackage{import}
\usepackage{xifthen}
\usepackage{transparent}
\usepackage[table]{xcolor}


\newcommand{\abbicon}{%
  \protect\tikz[baseline=0.2mm,scale=0.8]{%
    % Das blaue Quadrat mit abgerundeten Ecken
    \fill[uniblau, rounded corners=1.2pt] (0,0) rectangle (0.8em, 0.8em);
    % Der weiße Kreis in der Mitte
    \fill[white] (0.4em, 0.4em) circle (0.12em);
  }%
} 
\usepackage[labelfont={bf},font={color=miudiutex,small},figurename={\abbicon $~$Abbildung}]{caption}


\newcommand{\bckgrnd}{
\AddToHookNext{shipout/background}{
\begin{tikzpicture}[remember picture, overlay]
 \draw[draw=none,fill=uniblau!70, shift={(current page.north west)}] (-3,0) rectangle (23,-3.75);
 \draw[draw=none,fill=miugray, shift={(current page.north east)}] (-7,0) rectangle (3,-3.75);
\end{tikzpicture}
}
}



\newcommand{\incfig}[2]{%
\def\svgwidth{#2\columnwidth}
\import{C://LaTeX-Files/pngs/}{#1.pdf_tex}
}

\renewcommand{\textbf}[1]{\textsf{{\bfseries {#1}}}}



\titleformat{\section}[hang]
		{\LARGE \color{uniblau!75!black} \sffamily \bfseries}
		{\thesection}
		{6mm}
		{}
		[]
\titleformat{\subsection}[hang]
		{\large \sffamily \bfseries \color{black}}
		{\thesubsection}
		{4mm}
		{{\color{miugray}$\blacksquare ~$}}

\titleformat{\subsubsection}[block]
		{\normalsize \bfseries \sffamily \color{miudiutex}}%
		{\normalsize \color{black} \thesubsubsection}
		{2mm}
		{{\color{miugray}$\blacksquare ~$}}
		[ \color{black}]
		
\titleformat{\paragraph}
		{\normalsize \bfseries \sffamily \color{black}}
		{\footnotesize \theparagraph}
		{1mm}
		{}

\pagestyle{empty}
\titlespacing*{\section}{0pt}{20mm}{12mm}
\titlespacing{\section}{0pt}{20mm}{12mm}

\titlespacing*{\subsection}{0pt}{14mm}{4mm}
\titlespacing{\subsection}{0pt}{14mm}{4mm}

\titlespacing*{\subsubsection}{0pt}{6mm}{3mm}
\titlespacing{\subsubsection}{0pt}{6mm}{3mm}

\titlespacing*{\paragraph}{0pt}{6mm}{2mm}
\titlespacing{\paragraph}{0pt}{6mm}{2mm}


\let\oldmarginfigure\marginfigure
\let\endoldmarginfigure\endmarginfigure

\let\oldmarginfigure\marginfigure
\let\endoldmarginfigure\endmarginfigure

\renewenvironment{marginfigure}[1][0pt] % [1] steht für 1 Argument, [0pt] ist der Standardwert
  {\oldmarginfigure[#1]\captionsetup{labelfont={sf,bf,color=uniblau!75!black},font={color=miudiutex,small}}}
  {\endoldmarginfigure}
  
  
  
  
\renewcommand{\labelitemi}{\color{uniblau!70}$\bullet$}
\renewcommand{\labelitemii}{\color{uniblau!70}$\cdot$}
\newcommand{\plmtsquare}{{\color{uniblau!70}$\blacksquare~$}}
\newcommand{\splmt}[1]{{\color{uniblau!75!black} \sffamily #1}}
\newcommand{\fplmt}[1]{{\color{uniblau!75!black} \sffamily \bfseries #1}}

\newcommand{\antwort}[1]{
\begin{tcolorbox}[
	enhanced,
	enforce breakable,
	standard jigsaw,
	boxrule=0.7pt,
	colframe=black,
	sharp corners,
	boxsep=5pt,
	title=\bfseries \sffamily Antwort,
	opacityback=0,
	coltitle=miudiutex,
	attach boxed title to top text left={yshift=-\tcboxedtitleheight/2},
	boxed title style={colback=white,colframe=white},
	bottomrule at break=0pt,
	toprule at break=0pt,
	pad at break=16mm,
]
	\begin{center}
	#1
	\end{center}
	\end{tcolorbox}
}




\rowcolors{2}{miugray!50}{miugray!50}
\arrayrulecolor{white} % Weiße Gitterlinien
\setlength{\arrayrulewidth}{1.5pt} % Etwas dickere Linien
\renewcommand{\arraystretch}{1.4} % Mehr Zeilenhöhe für bessere Lesbarkeit}


%\setlist[itemize]{itemsep=0mm}
\setlist[enumerate]{itemsep=0mm}
\usepackage[table]{xcolor}
\usepackage{pgfplots}
    \pgfplotsset{
        every axis post/.style={
	yticklabel=\empty,
	xticklabel=\empty,
	y label style={at={(axis description cs:-0.2,1.1)}},
	x label style={at={(axis description cs:1.05,0.2)}},
	ticks=none,
	width=6cm,
	axis lines=middle,
%	xlabel near ticks,
        },
    }
\usetikzlibrary{positioning, arrows.meta, calc}

\usepackage{relsize}
\usepackage{cancel}

\newif\ifshowme
\showmetrue


\DeclareRobustCommand{\symPumpe}{%
    \tikz[baseline=-0.5ex, line width=0.8pt, x=1.2ex, y=1.2ex]{
        \draw[fill=white] (0,0) circle (1);
        \draw (0,1) -- (1,0);
        \draw (0,-1) -- (1,0);
    }% 
}

\DeclareRobustCommand{\symVentil}{%
    \tikz[baseline=-0.5ex, line width=0.8pt, x=1.2ex, y=1.2ex]{
        \draw[fill=white] (-0.5,-0.6) -- (0.5,0.6) -- (-0.5,0.6) -- (0.5,-0.6) -- cycle;
        \draw (0,0) -- (0.8,0);
        \draw (0.8,-0.3) -- (0.8,0.3);
    }%
}

\DeclareRobustCommand{\symBogen}{%
    \tikz[baseline=-0.5ex, line width=0.8pt, x=1.2ex, y=1.2ex]{
        \draw[rounded corners=0.6ex] (0, 1) |- (1, 0);
    }%
}

\begin{document}
%\newgeometry{top=1.8cm,bottom=4cm,left=1.5cm,right=7.5cm,marginparsep=-2.00cm,marginparwidth=4.75cm}
\newgeometry{top=3cm,bottom=3cm,left=4.5cm,right=4.5cm}

%\bckgrnd

\begin{center}
\LARGE
\color{uniblau}
\textbf{Übung 03} \par \vspace*{2mm}
\large \color{miugray!50!black} \textbf{Prozessbezogene Lebensmitteltechnologie}\par 
\normalsize Zusatzaufgaben
\end{center}
\par 
%\vspace*{6mm}


\color{miudiutex}



%\section*{Zusatzaufgaben}

\subsection*{Aufgabe 1}
%\tcbdocmarginnote[colback=unigelb!50,colframe=uniblau]{
%Wiederholungsaufgabe
%}
In \figurename{~\ref{masse}} (auf der nächsten Seite) sehen Sie den fünfstufigen Herstellungsprozess eines Fertigprodukts. 
\begin{itemize}
\item Bestimmen Sie den Massenstrom von Produkt $P$
\item Bestimmen Sie den Massenstrom von Rücklauf $A$
\end{itemize}


\begin{figure}[h]
\centering
\resizebox{.8\textwidth}{!}{
\begin{tikzpicture}[font=\sffamily, >={Latex[length=6pt, width=5pt]}]

    % --- Definition der Stile für die Blöcke ---
    \tikzset{
        block/.style={rectangle, draw=black, minimum width=1.6cm, minimum height=1.4cm, align=center},
		blaublock/.style={block, fill=uniblau!50},
        grayblock/.style={block, fill=miugray!50},
        redblock/.style={block, fill=miugray}
    }

    % --- Positionierung der Blöcke ---
    \node[blaublock] (I)   at (0, 0)      {I};
    \node[redblock]  (II)  at (0, -2.5)   {II};
    \node[redblock]  (III) at (-1.5, -5.0){III};
    \node[redblock]  (IV)  at (1.5, -5.0) {IV};
    \node[grayblock] (V)   at (1.5, -7.5) {V};

    % --- Systemgrenze (gestrichelte Linie) ---
    \draw[dashed, semithick, uniblau] (-3.2, 1) rectangle (3.2, -8.7);
    \node[above left, inner sep=2pt] at (3.2, 1.2) {\color{uniblau}Systemgrenze};

    % --- Stoffströme (Pfeile & Beschriftungen) ---
    
    % Feed F
    \draw[->] (-4.8, 0) -- (-2.7, 0);
    \draw[->] (-2.7, 0) -- (I.west);
    \node[left, align=left, inner sep=2pt] at (-5, 0) {Feed F \\ 1000 kg/h\\15\% TS\\85\% Wasser};

    % Stream D
    \draw[->] (I.east) -- (4.2, 0);
    \node[right, align=left, inner sep=2pt] at (4.5, 0) {D \\ 150 kg/h\\0\% TS};

    % Stream B
    \draw[->] (I.south) -- (II.north) node[midway, right] {B};

    % Stream C, E und G (Aufteilung)
    \draw[->]  (II.south) -- (0, -3.75) node[midway, right] {C};
    \draw[-]  (0, -3.75) -- (-1.5, -3.75) node[midway, below,xshift=5pt] {E 10\% TS};
    \draw[->] (-1.5, -3.75) -- (III.north);
    \draw[-]  (0, -3.75) -- (1.5, -3.75) node[right] {G};
    \draw[->] (1.5, -3.75) -- (IV.north);

    % Stream A (Rückführung)
    \draw[-]  (III.west) -- (-2.7, -5.0);
    \draw[->] (-2.7, -5.0) -- (-2.7, 0);
    \node[right, align=left] at (-2.7, -2.5) {A\\5\% TS};

    % Stream R (Rückführung)
    \draw[-]  (IV.east) -- (2.7, -5.0) -- (2.7, -2.5);
    \draw[->] (2.7, -2.5) -- (II.east);
    \node[left] at (2.7, -3.75) {R};

    % Stream L
    \draw[->] (IV.south) -- (V.north) node[midway, right] {L};

    % Stream K
    \draw[->] (III.south) -- (-1.5, -9.0) node[below, align=left] {K\\450 kg/h\\20\% TS};

    % Stream Product P
    \draw[->] (V.south) -- (1.5, -9.0) node[below, align=left] {Produkt P\\80\% TS};

    % Stream W
    \draw[->] (V.east) -- (4.2, -7.5) node[right, align=left] {W\\Wasser};

\end{tikzpicture}
}

\caption{Schematische Darstellung des Herstellungsprozesses. Abgeändert nach Singh und Heldman.}
\label{masse}
\end{figure}



\ifshowme
\antwort{
Produkt $P$:
\begin{align*}
F = D + W + K + P
\end{align*}
Für die Trockensubstanz gilt:
\begin{align*}
0,15 F &=  0 D + 0 W +0,2 K + 0,8 P \\
0,15 F &=  0 + 0 +0,2 \cdot 450 + 0,8 P \\
150 &= 90 + 0,8 P \\
P &= 75 ~\text{kg/h}
\end{align*}
Rücklauf $A$:
\begin{align*}
E &= A + K \\
0,1 E &= 0,05 A + 0,2 K \\
0,1 (A+K) &= 0,05 A + 90 \\ 
0,1 + 45 &= 0,05 A + 90 \\
0,05 A &= 45 \\
A &= 900 ~\text{kg/h}\\ \\
E &= 1.350 ~\text{kg/h}
\end{align*}
}
\else
\pagebreak
\fi

\subsection*{Aufgabe 2}
Stellen Sie die Formel für den Druckverlust durch Reibung (für laminare Strömung) nach dem Rohrdurchmesser $R$ um.

\ifshowme
\antwort{
\begin{align*}
\Delta p &= \frac{8\eta \cdot \Delta l}{\pi \cdot R^4} \cdot \dot{V} \\
R &= \sqrt[4]{\frac{8\eta \cdot \Delta l}{\pi \cdot \Delta p} \cdot \dot{V}}
\end{align*}
}
\pagebreak
\else
\par \vspace*{8cm}
\fi

\subsection*{Aufgabe 3}
Stellen Sie nun die Formel für das Geschwindigkeitsprofil der laminaren Strömung nach dem Rohrdurchmesser $R$ um

\ifshowme
\antwort{
\begin{align*}
v(r) &= -\frac{R^2}{4\eta} \cdot \frac{dp}{dx} \cdot \left[ 1 - \left(\frac{r}{R}\right)^2 \right] \\
R &= \sqrt{\frac{-4\eta \cdot v(r)}{\frac{dp}{dx}} + r^2}\\
R &= \sqrt{-4\eta \cdot v(r) \cdot \frac{dx}{dp} + r^2}
\end{align*}
}
\else
\pagebreak
\fi

\subsection*{Aufgabe 4}
%Für Lebensmittelchemiker spielt manchmal auch eine etwas überholte Druckeinheit eine Rolle. Früher hat man Gasdrücke häufig über Flüssigkeitssäulen gemessen. Dabei befindet sich die Flüssigkeit in einem U-förmigen Röhrchen, auf der einen Seite wird das zu messende Gas angeschlossen. Durch den Druckunterschied zwischen dem zu zu messenden Druck p und dem Luftdruck $p_L$ = 1.013 hPa sind beide Pegel unterschiedlich hoch. Da die eingefüllte
%Flüssigkeitsmenge bekannt ist, genügt eine Skala auf der Luftdruckseite, um den Druck zu bestimmen. Da man den Stand der Flüssigkeit misst, gibt man den Druck dann in Einheiten wie
%mmH2O oder mmHg an.
%
%\begin{itemize}
%
%\item Spielt der Durchmesser des Röhrchens für die Messung des Drucks über den
%Wasserstand eine Rolle?
%\item Gebt für die folgenden drei angelegten Drücke 
%\begin{itemize}
%\item $p_1$ = 1.013 hPa
%\item $p_2$ = 1.020 hPa
%\item $p_3$ = 990 hPa 
%\end{itemize}
%den Druckunterschied $\Delta p = p - p_L$ unter Verwendung von Wasser in der Einheit mmH2O und von Quecksilber ($\mathlarger{\rho}$ = 13,5 g/cm$^3$) in mmHg an.
%Wie hängt der Pegel rechts von der Dichte der verwendeten Flüssigkeit ab?
%\end{itemize}
In einem Wasserrohr fließen 60 L/h Wasser. Der Durchmesser der Zuleitung beträgt 20 mm, der Durchmesser an der Ableitung 10 mm. Der absolute Druck an der Zuleitung entspricht einem Umgebungsdruck von 1 bar. Die Ableitung liegt zu dem 0,5 m tiefer als die Zuleitung. Wie hoch ist der relative Druck und der absolute Druck im Ableitungsrohr? Welchen Einfluss hat der Höhenunterschied auf den Druck?
\ifshowme
\antwort{
Bernoulli:
\begin{align*}
p + \frac{1}{2} \rho v^2 + \rho g h = ~\text{const.}
\end{align*}
Geschwindigkeit $v$
\begin{align*}
v_1 &= \frac{\dot{V}}{A} = \frac{60~\text{L/h}}{\pi \cdot \left(\frac{0,02~\text{m}}{2}\right)^2} = \frac{1,6 \cdot 10^{-5}~\text{m}^3/\text{s}}{3,14 \cdot 10^{-4}} = 0,05~\text{m/s}\\ \\
v_2 &= 0,05 \cdot \frac{20^2}{10^2} = 0,2~\text{m/s}
\end{align*}
Bernoulli nach $p_2$ umstellen:
\begin{align*}
p_2 &= p_1 + \frac{1}{2} \rho v^2 + \rho g h - \frac{1}{2} \rho v^2 \\
&= 100.000~\text{Pa} + 0,5 \cdot 1000 \cdot (0,05^2 - 0,2^2) + 1000 \cdot 9,81 \cdot 0,5 \\
&= 104.886,25
\end{align*}
Relativer Druck = 4.886,25 Pa \pp 
Einfluss Höhenunterschied: \par 
Durch die Verengung würde normalerweise über den dynamischen Druck ein kleiner relativer Unterdruck entstehen, durch den vorhandenen Höhenunterschied ist der Druck an der Verengung aber größer als der Umgebungsdruck.
}
%\pagebreak
\else
\par \vspace*{8cm}
\fi




\subsection*{Aufgabe 5}
Berechnen Sie den Massenstrom von R. Bestimmen Sie anschließend den Wasseranteil, sowie den Anteil der Trockensubstanz in R

\begin{figure}[h]
\centering
\resizebox{.7\textwidth}{!}{
\begin{tikzpicture}[
    >={Latex[fill=white, length=3mm, width=2.5mm]}, % Offene Pfeilspitzen
    font=\sffamily, 
    line width=1pt
]

% --- Blöcke (Unit 1 und Unit 2) ---
\node[draw, thick, fill=uniblau!75, minimum width=3.5cm, minimum height=2.5cm] (Unit1) at (0,0) {};
\node[draw, thick, fill=miugray!50, minimum width=3.5cm, minimum height=2.5cm, right=2cm of Unit1] (Unit2) {};

% --- Systemgrenze (gestrichelte Linie) ---
% Die Koordinaten werden relativ zu den Blöcken berechnet, um alles passend einzurahmen
\draw[dashed, thick, color=uniblau, dash pattern=on 6pt off 4pt, rounded corners=2mm]
    ($(Unit1.north west) + (-0.8cm, 0.8cm)$) rectangle ($(Unit2.south east) + (0.8cm, -1.8cm)$);

% --- Feed (F) ---
\node[align=left, left] (F) at ($(Unit1.west) - (1.5cm, 0)$) {F \\ 100 kg/h \\ TS 10\%};
\draw[->] (F) -- (Unit1.west);

% --- Wasserstrom nach oben ---
\draw[->] (Unit1.north) -- ++(0, 1.5cm) node[above] {WASSER};

% --- Zwischenstrom (K) ---
\draw[->] (Unit1.east) -- (Unit2.west) node[midway, align=center, yshift=5mm] {K \\ TS 20\%};

% --- Produktstrom (P) ---
\draw[->] (Unit2.east) -- ++(1.5cm, 0) node[right, align=left] {P \\ TS 30\%};

% --- Rückstrom / Recycle (R) ---
% Der Pfad geht erst nach unten, dann horizontal zurück und wieder nach oben in Unit 1
\draw[->] (Unit2.south) -- ++(0, -1cm) coordinate (R_corner)
    -- (Unit1.south |- R_corner) node[midway, below, align=center,yshift=5mm] {R \\ 5\% TS}
    -- (Unit1.south);

\end{tikzpicture}
}
\caption{Modell zum Massenstrom}

\end{figure}

\ifshowme
\antwort{
\begin{align*}
    100  \cdot 0,1 &= P \cdot 0,3 \\
    10  &= P \cdot 0,3 \\
    P &= 33,33 ~\text{kg/h}
\end{align*}
Zweite Beziehung:
\begin{align*}
R + F &= W + K \\ 
0,05 R + 0,1F &= \cancel{W} + 0,2 K \\
0,05 R + 10 &= 0,2 K \\
0,05 R + 10 &= 0,2 (P + R) \\
0,05 R + 10 &= 6,66 + 0,2R \\
3,3333 &= 0,15R \\
R &= 22,22
\end{align*}
5\% TS von 22,22 kg/h Massenstrom sind 1,11 kg/h und 95\% Wasser sind 21,11 kg/h
}
\else
\fi



\end{document}
